\noindent
We investigated the magnetic response of several candidate transformer-core materials, including a CuNi alloy measured at two temperatures, by recording their $B$--$H$ curves under the same alternating magnetic field. A 400-turn primary coil was driven with the same current that magnetised a 500-turn secondary, and a resistor--integrator stage plus PicoScope converted the induced voltages into $H$ and $B$. Twenty steady-state loops were acquired for each sample, and the loop area together with the local $dB/dH$ gradients yielded estimates of the relative permeability and the power loss per unit volume.
$\\$
$\\$
The air control data are approximately linear, but the measured $\mu_{\mathrm{r}}=0.37\pm 0.03$ flags a large systematic calibration offset, so the absolute magnitudes should be treated with caution.
$\\$
$\\$
Among the ferromagnets, mild steel reached $\mu_{\mathrm{r,max}}\approx 1.2\times 10^{2}$ with a hysteresis loss of about $1.38\times 10^{6}\,\si{\watt\per\metre\cubed}$, while transformer iron achieved $\mu_{\mathrm{r,max}}\approx 3.3\times 10^{2}$ and only $4.73\times 10^{5}\,\si{\watt\per\metre\cubed}$ of loss. At \SI{10}{\celsius} the CuNi alloy still exhibits a small loop ($\mu_{\mathrm{r,max}}\approx 7.9$, loss $\approx 4.6\times 10^{3}\,\si{\watt\per\metre\cubed}$), but heating it above \SI{40}{\celsius} collapses the loop into a line with $\mu_{\mathrm{r}}=1.33\pm 0.01$ and almost no hysteresis, consistent with it crossing a Curie-like transition.
$\\$
$\\$
Within the shared calibration uncertainties, transformer iron therefore emerges as the preferred material, mild steel sacrifices efficiency through its broader loop, and CuNi becomes weakly paramagnetic near \SI{40}{\celsius}, rendering it unsuitable for transformer-core applications.
